\subsection{Reference fourth moments for terminal $(1,\allowbreak b,\allowbreak c)$, $4\le\allowbreak  b<\allowbreak c$}\label{proof:Q5-C178}
% Source Q5-C178: lines 21--731. Common decorations/assembly are invoked at the end.
\paragraph{Q178.0 The original local theorem and fixed direction}\label{part:Q178.0}
On a nonempty original finite probability space $(X,\allowbreak \mu)$, let $K$ be
real skew with $|K(x,\allowbreak y)|\le\allowbreak 1$. All expectations, inner products, spectra,
variances, covariances and row bounds use this same original $\mu$.
Put
\fitdisplay{
Jv=(\E v)\one,\quad B=-K,\quad A=J-B,\quad T=J+B=A^*,
\quad u=B\one,\quad p=1-u,\quad q=1+u,
}{}
\fitdisplay{
f_n=A^n\one,\quad\bar f_n=T^n\one,\quad h_n=Tf_{n-1},
\quad\delta_j=\|B^j\one\|_2^2,\quad d=\delta_1,
\quad\theta=13/32,\quad\kappa=19/32.
}{}
The receiver has a true outward arm of length $a\ge\allowbreak 3$.
The terminal residual is $(1,\allowbreak b,\allowbreak c)$, $4\le\allowbreak  b<\allowbreak c$, and $L$ is any
independently specified positive original connector length.
All connector edges point terminal-to-receiver. The terminal unit arm
points inward, its $b$-arm outward throughout, and its $c$-arm has first
edge inward and the rest outward. Thus $F=\allowbreak qf_bh_c$.
Genuine equal pairs at the two roots have fixed messages $f_\ell h_\ell$
or the complete pair reversal; their products $E_o,\allowbreak E_t$ lie in $[0,\allowbreak 1]$.
Every edge belongs to the original tree. For all these independent integers,
\fitdisplay{
\E[E_of_a\{T^L(E_tqf_bh_c)\}^4]\le1-d/3000.
\label{eq:Q178:T186.A}
}{Q178.T186.A}
The direction precedes the host, measure and kernel and stays fixed
across the analytic intervals. The associated four-terminal applications
use the stated receiving arm, rather than an arbitrary compound weight.

\paragraph{Q178.1 Finite tools and the true receiving weight}\label{part:Q178.1}
The finite framework gives $0\le\allowbreak  d\le\allowbreak 1/3$, $\|B\|^2\le\allowbreak \theta$, $|u|\le\allowbreak 1$,
and positive $A,\allowbreak T$ with kernels in $[0,\allowbreak 2]$.
For $n\ge\allowbreak 1$, $0\le\allowbreak  f_n,\allowbreak \bar f_n\le\allowbreak 2$ and their means are at most one;
for $n\ge\allowbreak 3$ the means are at most $1-2d$.
Equal pairs are in $[0,\allowbreak 1]$. Delete only zero-mass atoms; then
$p(x),\allowbreak q(x)\ge\allowbreak \mu(x)>\allowbreak 0$ and the true tree messages are strictly positive.
Let $G_n(z)=\allowbreak \sum_{j=\allowbreak 0}^{n-1}z^j$, $G_0=\allowbreak 0$, $\beta_1=\allowbreak \beta_2=\allowbreak 0$, and
$\beta_n=\allowbreak \sum_{j=\allowbreak 1}^{\lfloor(n-1)/2\rfloor}(-1)^{j-1}\delta_j$.
For $n\ge\allowbreak 3$, $\tau=\allowbreak d-\delta_2\le\allowbreak \beta_n\le\allowbreak  d$, and
\fitdisplay{
\one+G_n(B)u=\bar f_n+\sum_{j=3}^n\beta_j\bar f_{n-j}.
\label{eq:Q178:T186.1}
}{Q178.T186.1}
The spectral measure $\nu$ of $-B^2$ at $u$ has mass $d$, support
$[0,\allowbreak \theta]$, and $\int x^j\,\allowbreak d\nu=\allowbreak \delta_{j+1}$.
At $z=\allowbreak i\sqrt x$, for a finite symbol $A_0+zC_0$,
\fitdisplay{
\|(A_0+zC_0)(B)u\|_2^2=\int(A_0^2+xC_0^2)\,d\nu,
\qquad\E[(A_0+zC_0)(B)u]=-\int C_0\,d\nu.
\label{eq:Q178:T186.2}
}{Q178.T186.2}
Global Cauchy--Schwarz gives $d^2\le\allowbreak \delta_2$ and
$\delta_3\ge\allowbreak \delta_2^2/d$ for $d>\allowbreak 0$; these are not pointwise $x\ge\allowbreak  d$.
For $n\ge\allowbreak 2$, $G_n=\allowbreak U+zV$ has $1-x\le\allowbreak  U,\allowbreak V\le\allowbreak 1$;
$n=\allowbreak 1$ has $(U,\allowbreak V)=\allowbreak (1,\allowbreak 0)$.
For $n\ge\allowbreak 3$, $|G_n|\le\allowbreak 1$, while the true $G_2$ norm is
$\sqrt{1+\theta}<\allowbreak 6/5$.
Odd $n$ has squared modulus $(1+x^n)/(1+x)$; even $n=\allowbreak 2j$ has
$(1-(-x)^j)^2/(1+x)$, with odd $j\ge\allowbreak 3$ covered by $2x^2+x^5<\allowbreak 1$.
The actual nonnegative receiving upper is
\fitdisplay{
0\le f_a\le C_a=\one-G_a(-B)u-\beta_a\one,
\quad\E C_a=1-\beta_a-\beta_{a+1}\le1-2\tau\le1-19d/16.
\label{eq:Q178:T186.3}
}{Q178.T186.3}
The row bound, including the $a=\allowbreak 3$ predecessor $G_2$, gives
$C_a\le\allowbreak  p+(6/5)\sqrt d$.
For $Z=\allowbreak s\one+Q(B)u$ and $\zeta=\allowbreak -\E[Q(B)u]$,
\fitdisplay{
\operatorname{Var}Z=\|Q(B)u\|_2^2-\zeta^2,
\quad\operatorname{Cov}(C_a,Z)
=-\int\operatorname{Re}(G_aQ)\,d\nu-\beta_{a+1}\zeta.
\label{eq:Q178:T186.4}
}{Q178.T186.4}
The last index is $a+1$, and the original mean is retained.
For $W\ge\allowbreak 0$ and its actual mean $m=\allowbreak \E W$,
\fitdisplay{
W^4=m^4+4m^3(W-m)+(W-m)^2\{(W+m)^2+2m^2\}.
\label{eq:Q178:T186.5}
}{Q178.T186.5}
If $C_a\{(W+m)^2+2m^2\}\le\allowbreak  k$ pointwise, then
\fitdisplay{
\E[C_aW^4]\le(\E C_a)m^4+4m^3\operatorname{Cov}(C_a,W)
+k\operatorname{Var}W.
\label{eq:Q178:T186.6}
}{Q178.T186.6}
For $W=\allowbreak \alpha Z$ this is exactly
$\alpha^4[(\E C_a)(\E Z)^4+4(\E Z)^3\operatorname{Cov}(C_a,Z)]
+k\alpha^2\operatorname{Var}Z$.
Negative covariance uses a lower mean bound; positive covariance uses an
upper mean bound.
For a mean-zero original real vector $V$,
$|BV(x)|^2\le\allowbreak  pq\E V^2$. Thus for $Z=\allowbreak Cq+BV+c\one$,
$e=\allowbreak \E V$ and $\sigma^2=\allowbreak \operatorname{Var}V$,
\fitdisplay{
Z\le(C+e)q+(c-e)\one+\sqrt{pq}\,\sigma.
\label{eq:Q178:T186.7}
}{Q178.T186.7}
Only when $c=\allowbreak e$ has been proved does the free constant cancel;
this cancellation does not hold at $L=\allowbreak 1$.

\paragraph{Q178.2 The finite first-row packet and propagation}\label{part:Q178.2}
Use the actual envelope and signed first-row facts proved in parts C173.2--4
of Section~\ref{proof:Q5-C173}. This input concerns finite envelopes and
coefficients, rather than a conversion of its ordinary-fourth conclusion.
Let $n=\allowbreak c-b\ge\allowbreak 1$, $c_0=\allowbreak \beta_{c-1}-\beta_{b-1}$, and
\fitdisplay{
\Gamma=\beta_b+\beta_{b-1}+\beta_c+\beta_{c-1},
\quad r(z)=-c_0+(-1)^bz^bG_n(-z),\quad R=r(B)u.
}{}
Retain $\beta_b,\allowbreak \beta_{b-1}p$ from $f_b$ and $\beta_{c-1}q$ from $h_c$:
\fitdisplay{
f_b+h_c\le(2-\Gamma)\one+R,
\qquad F\le\alpha q e^R,\quad\alpha=e^{-\Gamma}.
\label{eq:Q178:T186.8}
}{Q178.T186.8}
This is AM--GM for the two actual nonnegative arms.
For $0<\allowbreak d\le\allowbreak 1/4$, $|R|<\allowbreak 13/50<\allowbreak 1/2$, and the Taylor bound
$e^R\le\allowbreak 1+R+(3/5)R^2$ gives
\fitdisplay{
T(qe^R)\le\bar f_2+BR+\delta\one,
\quad\delta=\E R+\E(uR)+\sqrt{d\E R^2}+(12/5)\E R^2.
\label{eq:Q178:T186.9}
}{Q178.T186.9}
The same $R$ and its signed means give the complete quarter-domain packet
\fitdisplay{
\begin{array}{c|c|c|c}
b&\Gamma/d\text{ lower}&D/d\quad(D\ge\delta_+)&|r|\text{ upper }\rho\\\hline
4&51/16&1&3/10\\
5&89/32&2/3&1/4\\
b\ge6&19/8&7/20&1/8.
\end{array}
}{}
Write $G_n(-z)=\allowbreak U-zV$. At $b=\allowbreak 4$,
$r=\allowbreak -c_0+x^2U-zx^2V$, $c_0\le\allowbreak 0$, $|c_0|\le\allowbreak \delta_2$.
At $b=\allowbreak 5$, $r=\allowbreak -c_0-x^3V-zx^2U$ with the same $c_0$ bounds.
For $b\ge\allowbreak 6$, $|c_0|\le\allowbreak \delta_3$; if $r=\allowbreak A_0+zC_0$ then
$|A_0|\le\allowbreak \delta_3+x^3$ and $|C_0|\le\allowbreak  x^3$.
These include $n=\allowbreak 1$ and both parities; $\rho$ is a finite-symbol bound,
not a pointwise quotient $R/u$.
For $D_i=\allowbreak \E B^iR$, the exact finite coefficients are
\fitdisplay{
D_{2j-1}=(-1)^j\int x^{j-1}A_0\,d\nu,
\qquad D_{2j}=(-1)^{j+1}\int x^jC_0\,d\nu.
\label{eq:Q178:T186.10}
}{Q178.T186.10}
At $b=\allowbreak 4$, $D_1,\allowbreak D_2\le\allowbreak 0$; at $b=\allowbreak 5$, $D_1\le\allowbreak \delta_4,\allowbreak D_2\le\allowbreak 0$.
For $b\ge\allowbreak 6$, $|D_1|<\allowbreak 67d/600$, $|D_2|<\allowbreak d/36$ and
$4(D_1+D_2)<\allowbreak 3d/5$.
Every $\gamma_j=\allowbreak \beta_j-D_{j-2}$, $j\ge\allowbreak 3$, satisfies
\fitdisplay{
\gamma_j\ge h:=15d/32.
\label{eq:Q178:T186.11}
}{Q178.T186.11}
Specifically $\gamma_3\ge\allowbreak  d/2$, $\gamma_4\ge\allowbreak 11d/12$;
the others use $|D_i|<\allowbreak d/8$ and $\tau\ge\allowbreak 19d/32$.
Positive propagation gives
\fitdisplay{
T^LF\le\alpha Y_L,\qquad
Y_L=\bar f_{L+1}+B^LR+D\bar f_{L-1}
+\sum_{i=1}^{L-1}D_i\bar f_{L-1-i}.
\label{eq:Q178:T186.12}
}{Q178.T186.12}
The first $Y_1=\allowbreak \bar f_2+BR+D\one$ is a nonnegative actual upper;
$Y_L=\allowbreak T^{L-1}Y_1\ge\allowbreak 0$. Finite convolution rewrites it as
\fitdisplay{
Y_L=\one+G_{L+1}(B)u+B^LR+D\bar f_{L-1}
-\sum_{j=3}^{L+1}\gamma_j\bar f_{L+1-j}.
\label{eq:Q178:T186.13}
}{Q178.T186.13}
At $L=\allowbreak 3,\allowbreak 4,\allowbreak \ge\allowbreak 5$ retain only the existing last two, three, four messages,
respectively. The bound $D\bar f_{L-1}\le\allowbreak  D[\one+G_{L-1}(B)u]$
uses this same finite convolution.

\paragraph{Q178.3 Finite symbol and expansion-coefficient lemmas}\label{part:Q178.3}
For $n\ge\allowbreak 2$, $G_n(-z)=\allowbreak U-zV$ satisfies
$0\le\allowbreak  U-(1-x)V\le\allowbreak  x$.
For even $n=\allowbreak 2j$, $U=\allowbreak V=\allowbreak a_j=\allowbreak (1-(-x)^j)/(1+x)$, giving difference
$xa_j$; for odd $n=\allowbreak 2j+1$ it is $(x+(-x)^j)/(1+x)$.
At $j=\allowbreak 1$ it is zero; even $j\ge\allowbreak 2$ gives at most $x$, and odd $j\ge\allowbreak 3$
lies between zero and $x$.
The $n=\allowbreak 1$ coordinates $(1,\allowbreak 0)$ are exceptional; the unit $(b,\allowbreak c)=\allowbreak (4,\allowbreak 5)$
is treated separately.
For receiving $G_a=\allowbreak E_a+zD_a$, even $a$ has $E_a=\allowbreak D_a$ and odd $a$
has $E_a=\allowbreak 1-xD_a$.
Thus $\operatorname{Re}(G_aG_n)\ge\allowbreak 0$ for $a\ge\allowbreak 3,\allowbreak n\ge\allowbreak 4$:
if either is even, use $(1-x)(1-2x)>\allowbreak 0$;
if both are odd, $n\ge\allowbreak 5$ has $D_n\le\allowbreak 1-x+x^2$, so the lower bound is
$1-3x+3x^2-3x^3+x^4\ge\psi(x)=1-3x+2x^2-x^3
\ge1451/32768>0$.
Here $\psi'=\allowbreak -3+4x-3x^2<\allowbreak 0$.
For $n=\allowbreak 3$ the uniform bound is only
\fitdisplay{
\operatorname{Re}(G_aG_3)\ge1-3x+x^2\ge-55/1024>-1/16.
\label{eq:Q178:T186.14}
}{Q178.T186.14}
For $n=\allowbreak 2$, $\operatorname{Re}(G_aG_2)\ge\allowbreak 0$.
Also $\operatorname{Re}[G_a(2+z)]\le\allowbreak 2$ and
$\operatorname{Re}[G_a(3+2z+z^2)]\le\allowbreak (3-x)E_a\le\allowbreak 3$.

\paragraph{Q178.3.2 The finite two-term shift}\label{part:Q178.3.2}
For $n\ge\allowbreak 3$, $0\le\allowbreak \sigma\le\allowbreak 3/16$,
$|G_n-\sigma(2+z)|\le\allowbreak 1-2\sigma$.
Set $A_0=\allowbreak 1-2\sigma$. At $n=\allowbreak 3,\allowbreak 4$ the squared differences from $A_0^2$,
divided by $x$, are respectively
$-1+2\sigma+\sigma^2+x$ and
$-1+2\sigma+\sigma^2+(-1+2\sigma)x+x^2$, strictly negative on the rectangle.
For $n\ge\allowbreak 5$, multiply by $1-z$ to obtain numerator $M-z^n$,
$M=\allowbreak A_0+\sigma z+\sigma z^2$, $|M|\le\allowbreak 1$, and
\fitdisplay{
A_0^2(1+x)-|M|^2=x[1-2\sigma-\sigma^2(1+x)]\ge(4715/8192)x.
}{}
The finite tail increment divided by $x$ is at most
$2\theta^{3/2}+\theta^4<\allowbreak 889/1600<\allowbreak 9/16<\allowbreak 4715/8192$.
At $x=\allowbreak 0$ it is direct. This is the finite proof also used by C102.

\paragraph{Q178.3.3 The finite three-term shift}
For $n\ge\allowbreak 4$, $0\le\allowbreak \sigma\le\allowbreak 49/288$,
$|G_n-\sigma(3+2z+z^2)|\le\allowbreak 1-3\sigma$.
Set $A_0=\allowbreak 1-3\sigma$. At $n=\allowbreak 4$ the quotient of the squared difference is
$-1+4\sigma-2\sigma^2+(-1+2\sigma+\sigma^2)x+x^2<\allowbreak 0$:
the constant is at most $-3/8$, linear coefficient negative, and $x^2<\allowbreak 1/6$.
At $n=\allowbreak 5$ it is
$-1+4\sigma-2\sigma^2+(1-4\sigma+\sigma^2)x+(-1+2\sigma)x^2+x^3$,
at most $-3/8+x-(95/144)x^2+x^3$.
The latter increases to $-3121/294912$ at $\theta$;
$1-4\sigma+2\sigma^2\ge\allowbreak 15649/41472>\allowbreak 3/8$.
For $n\ge\allowbreak 6$, multiplication by $1-z$ gives $M-z^n$,
$M=\allowbreak A_0+\sigma z+\sigma z^2+\sigma z^3$, and
\fitdisplay{
A_0^2(1+x)-|M|^2
=x[1-4\sigma+2\sigma^2+\sigma^2x-\sigma^2x^2]\ge(3/8)x.
}{}
Since $|z+z^2+z^3|^2=\allowbreak x-x^2+x^3\le\allowbreak 1$, $|M|\le\allowbreak  A_0+\sigma\le\allowbreak 1$.
The tail increment divided by $x$ is at most
$2\theta^2+\theta^5<\allowbreak 877/2560<\allowbreak 3/8$; $x=\allowbreak 0$ is direct.
Here $h\le\allowbreak 15/128<\allowbreak 49/288$, and every used index is at least four.

\paragraph{Q178.3.4 Continuous coefficient $17/2$ templates}\label{part:Q178.3.4}
Let $t=\allowbreak \sqrt d\le\allowbreak 1/2$, $W\ge\allowbreak 0$, and suppose
$\E W\le\allowbreak \beta\le\allowbreak  e^{-\gamma d}$,
$W+\E W\le\allowbreak \beta(q+1+r_0t+s_0t^2)$, and $C_a\le\allowbreak  p+(6/5)t$.
The pointwise coefficient is at most
\fitdisplay{
e^{-2\gamma t^2}(p+(6/5)t)
\{q^2+2q+3+2(q+1)(r_0t+s_0t^2)+(r_0t+s_0t^2)^2\}.
}{}
The constant is $6+q-q^3\le\allowbreak 13/2$, and the exact linear coefficient is
$(6/5-2r_0)q^2+(12/5+2r_0)q+18/5+4r_0$.
The next four coefficients have bounds
\fitdisplay{
2r_0^2+(9/2)s_0+(36/5)r_0,\quad
4r_0s_0+(6/5)r_0^2+(36/5)s_0,\quad
2s_0^2+(12/5)r_0s_0,\quad(6/5)s_0^2.
}{}
Compress higher degrees using $t\le\allowbreak 1/2$.
The linear bounds are obtained from the concave quadratic on $[0,\allowbreak 2]$:
\fitdisplay{
\begin{array}{c|c|c|c|c|c}
r_0&s_0&\gamma&\text{linear }\ell&\text{quadratic }B_*&17\gamma-B_*\\\hline
11/10&0&19/8&66/5&12&227/8\\
13/10&0&19/8&133/10&14&211/8\\
7/5&0&19/8&27/2&16&195/8\\
7/5&3/16&51/16&27/2&18&579/16\\
6/5&0&19/8&66/5&13&219/8\\
6/5&1&115/32&66/5&25&1155/32\\
13/10&1&145/48&133/10&26&1217/48\\
5/4&5/12&251/96&53/4&18&2539/96.
\end{array}
}{}
Every quadratic $2-\ell t+(17\gamma-B_*)t^2$ has negative discriminant;
compare with $(17/2)e^{2\gamma t^2}\ge\allowbreak 17/2+17\gamma t^2$ to obtain
coefficient at most $17/2$.
For example the third row has discriminant $-51/4$.
The case $r_0=\allowbreak 13/10,\allowbreak s_0=\allowbreak 0,\allowbreak \gamma=\allowbreak 89/32$ may use the weaker second row.

\paragraph{Q178.3.5 The unit large-$b$ coefficient eight}\label{part:Q178.3.5}
If $\E W\le\allowbreak \beta\le\allowbreak  e^{-(79/32)d}$ and
$W+\E W\le\allowbreak \beta[q+1+(23/20)t]$, the coefficient is at most eight.
The constant is at most $32/5$, since $2/(3\sqrt3)<\allowbreak 2/5$.
At $r_0=\allowbreak 23/20,\allowbreak s_0=\allowbreak 0$, the linear coefficient increases on $[0,\allowbreak 2]$
and is at most $66/5$; the compressed quadratic is
$(13/5)r_0^2+(36/5)r_0<\allowbreak 12$.
Compare with $8e^{(79/16)t^2}\ge\allowbreak 8+(79/2)t^2$:
the remaining quadratic is $8/5-(66/5)t+(55/2)t^2>\allowbreak 0$, with
discriminant $-44/25$. This coefficient is for the actual scaled $W$.

\paragraph{Q178.4 Every nonunit connector}\label{part:Q178.4}
Let $\alpha=\allowbreak e^{-\Gamma}$ and $h=\allowbreak 15d/32$.
The means $\mu_Z,\allowbreak \mu_Y$ refer to unscaled original uppers.
A symbol norm below is on the original spectral support and bounds
$\|P(B)u\|_2/\sqrt d$. The exponential payments use
$1-x\le\allowbreak  e^{-x}$ and actual nonnegative messages, with the direction fixed.

\paragraph{Q178.4.1 Every $L\ge\allowbreak 5$}
Retain $\bar f_0,\allowbreak \bar f_1,\allowbreak \bar f_2,\allowbreak \bar f_3$ and use
$D\bar f_{L-1}\le\allowbreak  DS'_{L-1}$:
\fitdisplay{
\begin{gathered}
Z=Cq+BP(B)u-h(1-d)\one,\quad C=1+D-3h\le1-13d/32,\\
P=G_L+z^{L-1}r+DG_{L-2}-h(2+z),\\
Q=G_{L+1}+z^Lr+DG_{L-1}-h(3+2z+z^2),\quad s=1+D-4h+hd.
\end{gathered}
}{}
Here $Z\ge\allowbreak  Y_L\ge\allowbreak 0$.
The finite shifts proved above give
$\|P\|\le\allowbreak 1-2h+D+\rho\theta^2<\allowbreak 11/10$ and
$\|Q\|\le\allowbreak  C+\rho\theta^{5/2}<\allowbreak C+1/30$.
For the first use the worst bound
$1+d/16+1/20\le\allowbreak 1+1/64+1/20<\allowbreak 11/10$;
for the second use $\sqrt\theta<\allowbreak 13/20$ and
$(3/10)\theta^{5/2}<\allowbreak 1/30$.
The true mean loss is
$\zeta=\allowbreak \gamma_{L+2}+D\beta_L-2hd>\allowbreak 0$, so
$\mu_Z\le\allowbreak 1+D-5h+3hd-D\tau\le\allowbreak 1-127d/128$.
The $G_aG_n$ products are nonnegative for $a\ge\allowbreak 3,\allowbreak n\ge\allowbreak 4$,
and the three-term shift has receiving real part at most three.
Therefore
$\operatorname{Cov}(C_a,Z)\le3hd+\rho\theta^{L/2}d
\le45d^2/32+d/30$,
and $\operatorname{Var}Z\le\allowbreak (C+1/30)^2d$.
The row $Z\le\allowbreak  q+(11/10)\sqrt d$, $\mu_Z\le\allowbreak 1$ satisfies the first
$17/2$ template for $W=\allowbreak \alpha Z$.
The first-term rate is at least
$4(19/8)+19/16+4(127/128)=\allowbreak 469/32$.
The variance term is at most
$(16337/1800)d e^{-(1373/248)d}\le\allowbreak (16337/1800)de^{-11d/2}$;
the covariance term is at most
$[(2/15)d+(45/8)d^2]e^{-(1597/128)d}$.
Relative to $de^{-11d/2}$ its extra coefficient is at most
$2/15+270/893$, and $16337/1800+2/15+270/893<\allowbreak 10$.
Thus this entire row is bounded by $e^{-469d/32}+10de^{-11d/2}$,
which is dominated by the common scalar upper below.

\paragraph{Q178.4.2 $L=\allowbreak 4$}
Retain only $\bar f_0,\allowbreak \bar f_1,\allowbreak \bar f_2$:
\fitdisplay{
\begin{gathered}
C=1+D-2h,\quad s=1+D-3h,\\
Q=G_5+z^4r+DG_3-h(2+z),\quad
P=G_4+z^3r+DG_2-h,\quad Z=Cq+BP(B)u-h\one.
\end{gathered}
}{}
The loss $\zeta=\allowbreak \gamma_6+(D-h)d>\allowbreak 0$ gives
$\mu_Z\le\allowbreak 1+D-4h-(D-h)d\le\allowbreak 1-7d/8$.
Also $Cq-h\le\allowbreak  q-11d/32$ and
$\|P\|\le\allowbreak 1-h+(6/5)D+\rho\theta^{3/2}<\allowbreak 13/10$.
Indeed $|G_4-h|\le\allowbreak 1-h$, since its squared difference quotient is
$-1+2h-x+x^2<\allowbreak 0$.
The second template gives $17/2$.
Moreover $\|Q\|\le\allowbreak 1-2h+D+\rho\theta^2\le\allowbreak 21/20+d/16$.
The actual $DG_3$ exception costs $Dd/16$, so
$\operatorname{Cov}(C_a,Z)\le2hd+(D/16)d+\rho\theta^2d
\le d^2+d/20$.
The three fourth-expansion terms are at most
\fitdisplay{
e^{-227d/16}+(7497/800)de^{-389d/84}
+[(1/5)d+4d^2]e^{-97d/8}.
}{}
Relative to $de^{-9d/2}$ the coefficient is at most
$7497/800+1/5+12/61<\allowbreak 10$.

\paragraph{Q178.4.3 $L=\allowbreak 3$, with all three $b$ ranges}
Here
\fitdisplay{
\begin{gathered}
Z=Cq+BP(B)u-h\one,\quad C=1+D-h,\quad s=1+D-2h,\\
Q=G_4+z^3r+DG_2-h,\quad P=G_3+z^2r+D.
\end{gathered}
}{}
The positive loss is $\zeta=\allowbreak \gamma_5+Dd$, and
$\mu_Z\le\allowbreak 1+D-3h-Dd$.
Also $\operatorname{Cov}(C_a,\allowbreak Z)\le\allowbreak  hd+\rho\theta^{3/2}d$ and
$\|Q\|\le\allowbreak 1-h+(6/5)D+\rho\theta^{3/2}$, using $|G_4-h|\le\allowbreak 1-h$.
The pointwise row is
$Z+\mu_Z\le\allowbreak  q+1+\|P\|\sqrt d+2(D-h)_++D-4h$.
If $D\ge\allowbreak  h$, its last constant is $3D-6h$;
if $D<\allowbreak h$, retain $D-4h$ rather than setting $q=\allowbreak 2$ against a negative coefficient.
For $b=\allowbreak 4$, $\mu_Z\le\allowbreak 1-13d/32$, $\|P\|<\allowbreak 7/5$, and
$3D-6h=\allowbreak 3d/16$. The fourth template gives $17/2$.
Also $\|Q\|\le\allowbreak 27/25+(117/160)d$.
The variance term is at most $(12393/1250)de^{-5d}$, with
$12393/1250=\allowbreak (17/2)(27/25)^2$.
The first rate is $249/16>\allowbreak 31/2$.
The covariance has quadratic cost $(15/8)d^2$, rate $447/32$,
and gap $287/32$ relative to $de^{-5d}$.
Using $e>\allowbreak 8/3$ adds $45/574$; the total
$12393/1250+8/25+45/574<\allowbreak 52/5$.
Thus the upper is $e^{-31d/2}+(52/5)de^{-5d}$.
It is at most $F_0=\allowbreak e^{-25d/2}+(103/10)de^{-9d/2}$:
the extra first rate of three provides at least $3de^{-31d/2}$,
whereas the extra coefficient costs at most $de^{-5d}/10$;
$e^{(21/2)d}<\allowbreak e^3<\allowbreak 27<\allowbreak 30$ on the quarter-domain pays it.
For $b=\allowbreak 5$, $\mu_Z\le\allowbreak 1-71d/96$, $\|P\|<\allowbreak 13/10$, and $3D-6h<\allowbreak 0$.
Use the second template with its stronger $\gamma$.
One has $\|Q\|\le\allowbreak 16/15+(53/160)d$, variance coefficient $2176/225$,
and rate $1265/256>\allowbreak 9/2$.
The first rate is $733/48>\allowbreak 25/2$; the covariance linear coefficient is
$4/15$, with normalized quadratic cost at most $45/566$.
Since $2176/225+4/15+45/566<\allowbreak 103/10$, this is at most $F_0$.
For $b\ge\allowbreak 6$, $\mu_Z\le\allowbreak 1-169d/160$, $\|P\|<\allowbreak 6/5$,
$D<\allowbreak h$ and $D-4h<\allowbreak 0$. The fifth template gives $17/2$.
Here $\|Q\|\le\allowbreak 31/30-(39/800)d$, variance coefficient $16337/1800$
with rate greater than $9/2$, and first rate $1193/80>\allowbreak 25/2$.
The covariance linear coefficient is $2/15$, with normalized quadratic
cost at most $225/2614$; the total is less than $103/10$.

\paragraph{Q178.4.4 $L=\allowbreak 2$: exact constant cancellation}
The actual form is
$Y_2=\allowbreak Cq+BP(B)u+c\one$, with
$C=\allowbreak 1+D$, $P=\allowbreak G_2+zr$, and $c=\allowbreak D_1-d=\allowbreak \E(P(B)u)$.
The centered row yields
$Y_2\le\allowbreak (C+c)q+\sqrt{pq\operatorname{Var}(P(B)u)}$,
where $C+c=\allowbreak 1+D-d+D_1\le\allowbreak 1$.
Also $\|P\|\le\allowbreak 6/5+\rho\sqrt\theta<\allowbreak 7/5$ and $\mu_{Y_2}\le\allowbreak 1$.
The third template gives $17/2$ uniformly.
The full symbol is $Q=\allowbreak G_3+D+z^2r$, with $\zeta=\allowbreak d-D_2>\allowbreak 0$.
Its real and imaginary coordinates are positive:
$|A_0|\le\allowbreak \delta_2+x^2$, $|C_0|\le\allowbreak  x^2$, and
$D\ge\allowbreak 7d/20>\allowbreak \theta^2d$ give real part at least
$1-x-x^3+D-x\delta_2>\allowbreak 0$ and imaginary coordinate at least $1-x^3>\allowbreak 0$.
Thus their positive uppers may be squared.
For $b=\allowbreak 4$, the real coordinate is at most $1+D-x$ since $A_0\ge\allowbreak 0$,
and the imaginary coordinate at most $1+x^3$.
The squared difference quotient is at most
$-1-2D+x+2x^3+x^6<\allowbreak 0$, so $\|Q\|\le\allowbreak 1+D$.
Then $\operatorname{Var}Y_2\le\allowbreak (1+d)^2d-d^2\le\allowbreak  de^{5d/4}$.
The mean is at most $1-d$, with first rate $287/16$ and variance
rate $41/8>\allowbreak 9/2$.
The covariance is at most $d(1/16+39/320)$, giving fourth cost
$59/80$. The combined coefficient $17/2+59/80<\allowbreak 103/10$ is dominated
by $F_0$.
For $b=\allowbreak 5$, real coordinate is at most $1+D-x+x^4$, imaginary at most
$1+x^3$. The quotient is at most $-1+x+4x^3+x^6+x^7<\allowbreak 0$, all $D$
terms being nonpositive; hence $\|Q\|\le\allowbreak 1+D$.
The variance is at most $(1+2d/3)^2d-d^2\le\allowbreak  de^{4d/9}$,
and the mean at most $1-5d/4$.
The first rate is $277/16>\allowbreak 25/2$, the variance rate
$89/16-4/9>\allowbreak 9/2$.
The covariance bound $d(1/16+13/128)=\allowbreak 21d/128$ gives cost $21/32$;
$17/2+21/32<\allowbreak 103/10$.
For $b\ge\allowbreak 6$ use $|A_0|\le\allowbreak \delta_3+x^3$, $|C_0|\le\allowbreak  x^3$ and
$\widetilde D=\allowbreak D+\theta^3d<\allowbreak 21d/50$.
The real coordinate is at most $1+\widetilde D-x+x^4$, imaginary at most
$1+x^4\le\allowbreak 1+x^3$, so $\|Q\|\le\allowbreak 1+\widetilde D$ as above.
Since $\zeta\ge\allowbreak 35d/36$, the variance is at most $d$:
$21/25-(35/36)^2$ plus the worst cubic contribution $441/10000$
is still negative.
The mean is at most $1-3d/2$, with first rate $267/16>\allowbreak 25/2$ and
variance rate $19/4>\allowbreak 9/2$.
The covariance bound $d(1/16+13/256)=\allowbreak 29d/256$ gives cost $29/64$;
$17/2+29/64<\allowbreak 103/10$.
In these cases use
$\operatorname{Re}(G_aG_3)\ge\allowbreak 1-3x+x^2\ge\allowbreak -55/1024>\allowbreak -1/16$;
the $DG_a$ term is nonnegative, the $r$ remainder costs $\rho\theta$,
and the mean correction $-\beta_{a+1}\zeta$ is nonpositive.

\paragraph{Q178.4.5 A common continuous scalar payment}
For $F_0=\allowbreak e^{-25d/2}+(103/10)de^{-9d/2}$, degree-four Taylor bounds give
\fitdisplay{
\begin{aligned}
\frac{e^{9d/2}-e^{-8d}}d
&\ge25/2-(175/8)d+(4825/48)d^2-(58975/384)d^3\\
&\ge25/2-(175/8)d+62d^2\ge25/2-30625/15872>211/20.
\end{aligned}
}{}
Since $211/20=\allowbreak 103/10+1/4$ and $e^{9/8}<\allowbreak 4$, this proves
$F_0\le\allowbreak 1-d/16$ for $0<\allowbreak d\le\allowbreak 1/4$.
For the latter exponential bound use $e<\allowbreak 3$, from a geometric factorial
tail, and $e^{1/8}\le\allowbreak 8/7$, so $e^{9/8}<\allowbreak 24/7<\allowbreak 4$.
Every nonunit row above is dominated by $F_0$, completing all $L\ge\allowbreak 2$
on the low domain.

\paragraph{Q178.5 Every unit connector}\label{part:Q178.5}
For $L=\allowbreak 1$, $Y_1=\allowbreak q+B(u+R)+D\one$.
Write $r=\allowbreak A_0+zC_0$, $P=\allowbreak 1+r$, $Q=\allowbreak 1+zP$, and
\fitdisplay{
e=\E R=-\int C_0\,d\nu,\quad
\zeta=-\E[Q(B)u]=\int(1+A_0)\,d\nu,
\quad\mu=\E Y_1=1-d+D+D_1.
}{}
The original centered row gives
$Y_1\le\allowbreak (1+e)q+(D-e)\one+\sqrt{pq}\sigma$,
where $\sigma^2=\allowbreak \operatorname{Var}((1+r)(B)u)$.
Set the actual upper $W=\allowbreak \alpha Y_1$.
For its coefficient bound use $\beta=\allowbreak \alpha(1+e)>\allowbreak 0$, retaining
$\alpha$ as the true scale in the fourth expansion.
The variance and covariance of $Y_1$ equal those of $Q(B)u$.
The $V$ in the finite coordinates $G_n(-z)=\allowbreak U-zV$ is a spectral scalar,
and the receiving mean correction remains $-\beta_{a+1}\zeta$.

\paragraph{Q178.5.1 $b=\allowbreak 4,\allowbreak c=\allowbreak 5$}
Here $n=\allowbreak 1$, $c_0=\allowbreak \beta_4-\beta_3=\allowbreak 0$, $R=\allowbreak x^2u$,
$(U,\allowbreak V)=\allowbreak (1,\allowbreak 0)$, and $e=\allowbreak 0$.
One has $\Gamma=\allowbreak 4d-\delta_2\ge\allowbreak 115d/32$, $D=\allowbreak d$,
$\mu=\allowbreak 1-\delta_3\le\allowbreak 1$, $Q=\allowbreak 1+z(1+x^2)$ and $\zeta\ge\allowbreak  d$.
Thus $|Q|^2=\allowbreak 1+x+2x^3+x^5\le\allowbreak 25/16$ and
$\operatorname{Var}Y_1\le\allowbreak 25d/16-d^2$.
The product has $\operatorname{Re}(G_aQ)\ge\allowbreak 1-2x-x^3>\allowbreak 0$, so its
receiving covariance is nonpositive.
Since $\|1+r\|\le\allowbreak 1+\theta^2<\allowbreak 6/5$,
$Y_1+\mu\le\allowbreak  q+1+(6/5)\sqrt d+d$.
The template $r_0=\allowbreak 6/5,\allowbreak s_0=\allowbreak 1,\allowbreak \gamma=\allowbreak 115/32$ has linear coefficient
at most $66/5$ and merged quadratic less than $25$.
The successive contributions $801/50,\allowbreak 858/125,\allowbreak 61/50,\allowbreak 3/20$ sum to
less than $25$. The remaining quadratic
$2-(66/5)t+(1155/32)t^2>\allowbreak 0$ has negative discriminant, proving the
coefficient throughout $0\le\allowbreak  t\le\allowbreak 1/2$.
Put
$H=\allowbreak 1+(115/16)d+(13225/512)d^2\le\allowbreak  e^{2\Gamma}$,
$A=\allowbreak 1-19d/16$, and $V_*=\allowbreak 25d/16-d^2$.
The reference fourth is at most $A/H^2+(17/2)V_*/H$.
For $G=\allowbreak H^2-A-(17/2)V_*H$, the coefficients at degrees one through
three are $73/32,\allowbreak 8377/512,\allowbreak 1463835/16384$; degree four is
$b_H^2+(17/2)b_H>\allowbreak 0$, $b_H=\allowbreak 13225/512$.
Hence $G\ge\allowbreak 73d/32$. Since $H(1/4)<\allowbreak 5$,
$G/H^2\ge\allowbreak 73d/800>\allowbreak d/16$, proving the unit bound $1-d/16$.

\paragraph{Q178.5.2 $b=\allowbreak 4,\allowbreak c\ge\allowbreak 6$}
For $n\ge\allowbreak 2$,
$r=\allowbreak -c_0+x^2U-zx^2V$, $c_0\le\allowbreak 0$, $|c_0|\le\allowbreak \delta_2$.
The finite relation $U-(1-x)V\le\allowbreak  x$ was proved in part~\ref{part:Q178.3}.
Both coordinates of $Q$ are positive:
$Q_{\rm re}=\allowbreak 1+x^3V$ and $Q_{\rm im}=\allowbreak 1-c_0+x^2U$,
where the latter denotes the coefficient of $z$.
For odd receiving index, $D_a=\allowbreak 1$ is the lower extremum, giving
$\operatorname{Re}(G_aQ)\ge\allowbreak 1-2x+c_0x-x^4$.
The even-index value is at least this, with difference containing
$x^2Q_{\rm im}\ge\allowbreak 0$.
Let $v=\allowbreak \delta_2/d\in[d,\allowbreak \theta]$. Then
$\Gamma\ge\allowbreak 4d-2\delta_2$ and
$\mu=\allowbreak 1+D_1\le\allowbreak 1-\kappa\delta_3\le\allowbreak 1-\kappa v^2d=\allowbreak :M(v)$.
Also $\mu\ge\allowbreak 1-\delta_2d-\delta_3\ge\allowbreak 1-(273/1024)d\ge\allowbreak 1-4d/15$,
and $\zeta\ge\allowbreak  d$.
Keep the true mean correction to get
\fitdisplay{
\begin{aligned}
\operatorname{Cov}(C_a,Y_1)
&\le-d+2\delta_2+\delta_2^2+\delta_5-\tau d\\
&\le-(1-2\theta-\theta^4)d-(1-\theta-\theta^2)d^2\\
&\le-4d/25-2d^2/5.
\end{aligned}
}{}
Here $1-2\theta-\theta^4=\allowbreak 168047/1048576>\allowbreak 4/25$ and
$1-\theta-\theta^2=\allowbreak 439/1024>\allowbreak 2/5$.
Since $\mu^3\ge\allowbreak 1-4d/5$,
\fitdisplay{
4\mu^3\operatorname{Cov}(C_a,Y_1)\le-L(d),
\quad L(d)=(16d/25+8d^2/5)(1-4d/5)
=16d/25+136d^2/125-32d^3/25.
}{}
The signed norm satisfies
\fitdisplay{
|Q|^2\le1+x+4x^3+x^5+x^6+2\delta_2x(1+x^2)+\delta_2^2x.
}{}
Subtracting $\zeta^2\ge\allowbreak  d^2$ gives
$\operatorname{Var}Y_1\le\allowbreak  V_*:=\allowbreak 17d/10-3d^2/5+7d^3/100$,
using $1+\theta+4\theta^3+\theta^5+\theta^6<\allowbreak 17/10$,
$2\theta^2(1+\theta^2)-1<\allowbreak -3/5$ and $\theta^3<\allowbreak 7/100$.
For the coefficient, $0\le\allowbreak  e\le\allowbreak \delta_3\le\allowbreak  d/6$,
$\|1+r\|\le\allowbreak 13/10$, and
$\beta=\allowbreak \alpha(1+e)\le\allowbreak  e^{-(145/48)d}$.
Because $\mu\le\allowbreak 1\le\allowbreak 1+e$,
$W+\E W\le\allowbreak \beta[q+1+(13/10)\sqrt d+d]$.
The template $r_0=\allowbreak 13/10,\allowbreak s_0=\allowbreak 1,\allowbreak \gamma=\allowbreak 145/48$ gives linear upper
$133/10$, merged quadratic less than $26$;
$2-(133/10)t+(1217/48)t^2>\allowbreak 0$ has negative discriminant.
For the same-$v$ endpoint replacement define
$A(d,\allowbreak v)=\allowbreak [1-2(1-v)d]M(v)^4-L(d)$.
This is positive: $M\ge\allowbreak 39/40$, $M^4\ge\allowbreak 9/10$,
$E_0=\allowbreak 1-2(1-v)d\ge\allowbreak 1/2$ and $L\le\allowbreak 13/50$, giving $A\ge\allowbreak 19/100$.
The derivative of $e^{8vd}A$, divided by $e^{8vd}d$, is
\fitdisplay{
2M^4+8E_0M^3(M-\kappa v)-8L>0.
}{}
Indeed, $2M^4\ge\allowbreak 9/5$, $E_0\ge\allowbreak 1/2$, $M^3\ge\allowbreak 7/10$,
$M-\kappa v\ge\allowbreak 7/10$ and $8L\le\allowbreak 52/25$, giving $42/25>\allowbreak 0$.
Consequently $e^{-4\Gamma}A(d,\allowbreak v)\le\allowbreak  e^{-(51/4)d}A(d,\allowbreak \theta)$.
Since $\kappa\theta^2>\allowbreak 1/12$, $M(\theta)\le\allowbreak 1-d/12$;
using $(1-d/12)^4\le\allowbreak 1-d/3+d^2/24$ gives
\fitdisplay{
A(d,\theta)\le\overline A
=1-(2593/1200)d-(1301/2000)d^2+(11813/9600)d^3.
}{}
Put $H=\allowbreak 1+(51/8)d+(2601/128)d^2\le\allowbreak  e^{2\Gamma_{\min}}$.
The moment is at most $\overline A/H^2+(17/2)V_*/H$.
For $G=H^2-\overline A-(17/2)V_*H
=d[g_1+g_2d+g_3d^2+g_4d^3+g_5d^4]$,
\fitdisplay{
\begin{gathered}
g_1=553/1200>9/20,\quad g_2=-5087/1000>-51/10,
\quad g_3=-14813/3840>-4,\\
g_4=6765201/16384+638979/6400,\qquad g_5=-309519/25600.
\end{gathered}
}{}
For $d\le\allowbreak 1/4$, $g_4+g_5d>\allowbreak 509$.
The bracket is at least $P(d)=\allowbreak 9/20-51d/10-4d^2+509d^3$.
For $d\le\allowbreak 1/16$, $P\ge\allowbreak 37/320>\allowbreak 1/10$.
For $d\ge\allowbreak 1/16$, use $509d^3\ge\allowbreak 509d^2/16$ to get
$P\ge\allowbreak 9/20-51d/10+(445/16)d^2\ge\allowbreak 9/20-2601/11125>\allowbreak 1/10$.
Thus $G\ge\allowbreak  d/10$. Since $H(1/4)=\allowbreak 7913/2048<\allowbreak 4$, the reference fourth
is at most $1-d/160$.

\paragraph{Q178.5.3 $b=\allowbreak 5,\allowbreak c\ge\allowbreak 6$}
The finite coefficients are
$A_0=\allowbreak -c_0-x^3V$, $C_0=\allowbreak -x^2U$, $c_0\le\allowbreak 0$, $|c_0|\le\allowbreak \delta_2$.
Here $\Gamma\ge\allowbreak 89d/32$, $D=\allowbreak 2d/3$, $0\le\allowbreak  e\le\allowbreak \delta_3\le\allowbreak  d/6$,
and $\mu=\allowbreak 1-d/3+D_1\le\allowbreak 1-d/4$, since
$D_1\le\allowbreak \delta_4\le\allowbreak \theta^3d<\allowbreak d/12$.
Also $\mu\ge\allowbreak 1-d/3-\delta_2d\ge\allowbreak 1-7d/16$ and
$\zeta=\allowbreak \int(1-c_0-x^3V)\,\allowbreak d\nu\ge\allowbreak (1-\theta^3)d\ge\allowbreak 11d/12$.
The coefficient scale satisfies
$\beta=\allowbreak \alpha(1+e)\le\allowbreak  e^{-(251/96)d}$, and
\fitdisplay{
W+\E W\le\beta[q+1+(5/4)\sqrt d+(5/12)d].
}{}
Here $\|1+r\|\le\allowbreak 5/4$ and the remaining constant
$D-d/4-2e\le\allowbreak 5d/12$.
The corresponding template gives linear upper $53/4$ and merged
quadratic less than $18$; the contributions are
$14,\allowbreak 167/48,\allowbreak 115/288,\allowbreak 5/192$.
The quadratic $2-(53/4)t+(2539/96)t^2>\allowbreak 0$ has negative discriminant.
The coordinates $Q_{\rm re}=\allowbreak 1+x^3U$ and $Q_{\rm im}=\allowbreak 1-c_0-x^3V$
are positive and
$|Q|^2\le\allowbreak 1+x+2x^3+x^6+2\delta_2x+\delta_2^2x$.
After the exact negative mean square,
$\operatorname{Var}Y_1\le\allowbreak  V_*:=\allowbreak 25d/16-d^2/2+7d^3/100$,
using $1+\theta+2\theta^3+\theta^6<\allowbreak 25/16$,
$2\theta^2-(11/12)^2<\allowbreak -1/2$ and $\theta^3<\allowbreak 7/100$.
Since $\operatorname{Re}(G_aQ)\ge\allowbreak 1-2x-\delta_2x$,
\fitdisplay{
\operatorname{Cov}(C_a,Y_1)
\le-3d/16+[\theta^2-\kappa(11/12)]d^2
\le-3d/16-3d^2/8,
}{}
where $\kappa(11/12)-\theta^2=\allowbreak 1165/3072>\allowbreak 3/8$.
Use $\mu^3\ge\allowbreak 1-21d/16$, $\E C_a\le\allowbreak 1-19d/16$ and
$\mu^4\le\allowbreak (1-d/4)^4\le\allowbreak 1-d+3d^2/8$.
The mean-plus-covariance bracket is at most
$\overline A=\allowbreak 1-47d/16+67d^2/64+195d^3/128$.
Let $H=\allowbreak 1+(89/16)d+(7921/512)d^2$.
For $G=H^2-\overline A-(17/2)V_*H
=d[g_1+g_2d+g_3d^2+g_4d^3+g_5d^4]$,
\fitdisplay{
\begin{aligned}
g_1&=25/32>3/4,\qquad g_2=-4501/512>-9,\\
g_3&=a_H^3-195/128-(17/2)[(25/16)b_H-a_H/2+7/100]>-12,\\
g_4&=b_H^2+(17/2)(b_H/2-(7/100)a_H),\\
g_5&=-(119/200)b_H=-942599/102400,
\qquad a_H=89/16,\quad b_H=7921/512.
\end{aligned}
}{}
On $d\le\allowbreak 1/4$, $g_4+g_5d>\allowbreak 299$, so the bracket is at least
$P(d)=\allowbreak 3/4-9d-12d^2+299d^3$.
For $d\le\allowbreak 1/16$, $P\ge\allowbreak 9/64$.
For $d\ge\allowbreak 1/16$, $P''=\allowbreak -24+1794d>\allowbreak 80$.
At $d_0=\allowbreak 1/9$, $P(d_0)=\allowbreak 35/2916$ and $P'(d_0)=\allowbreak -16/27$.
The tangent quadratic gives
$P(d)\ge\allowbreak 35/2916-8/3645=\allowbreak 143/14580>\allowbreak 1/200$.
Thus $G\ge\allowbreak  d/200$. Since $H(1/4)=\allowbreak 27505/8192<\allowbreak 7/2$, $H^2<\allowbreak 13$,
the reference fourth is at most $1-d/2600$, hence at most $1-d/3000$.

\paragraph{Q178.5.4 Every $b\ge\allowbreak 6,\allowbreak c>\allowbreak b$}\label{part:Q178.5.4}
The actual alternating moments give
$\beta_b+\beta_{b-1}\ge\allowbreak 2\tau$ and
$\beta_c+\beta_{c-1}\ge\allowbreak 2\tau+\delta_3$ for $c\ge\allowbreak 7$.
At $c=\allowbreak 7$ equality holds; for $c\ge\allowbreak 8$ each $\beta$ is at least
$\tau+\kappa\delta_3$, with $2\kappa>\allowbreak 1$.
For $v=\allowbreak \delta_2/d$, the global inequality $\delta_3\ge\allowbreak  v^2d$ yields
\fitdisplay{
\Gamma\ge4\tau+\delta_3\ge(4-4v+v^2)d
\ge(2601/1024)d>5d/2.
}{}
This imposes no pointwise lower spectral bound.
Also $|e|\le\allowbreak \delta_4\le\allowbreak \theta^3d<\allowbreak 7d/100$.
The actual $D_1<\allowbreak 67d/600$ gives $\mu\le\allowbreak 1-d/2$, with loss
$323/600>\allowbreak 1/2$; and
$\mu\ge\allowbreak 1-13d/20-\delta_3d-\delta_4\ge\allowbreak 1-4d/5$,
$\zeta\ge\allowbreak  d-\delta_3d-\delta_4\ge\allowbreak 7d/8$.
The scale $\beta=\allowbreak \alpha(1+e)>\allowbreak 0$ obeys $\beta\le\allowbreak  e^{-(79/32)d}$,
since $2601/1024-\theta^3=\allowbreak 81035/32768>\allowbreak 79/32$.
The centered noise has bound
\fitdisplay{
\frac{\sigma}{(1+e)\sqrt d}
\le\frac{9/8}{1-\theta^3/4}
<\frac{9/8}{393/400}=450/393<23/20.
}{}
Moreover $\mu\le\allowbreak 1-d/2\le\allowbreak 1+e$ and
$D-d/2-2e\le\allowbreak 7d/20-d/2+14d/100=\allowbreak -d/100<\allowbreak 0$.
Thus $W+\E W\le\allowbreak \beta[q+1+(23/20)\sqrt d]$.
For coefficient eight, use
$p(q^2+2q+3)=\allowbreak 6+q-q^3\le\allowbreak 32/5$, since $2/(3\sqrt3)<\allowbreak 2/5$.
At $r_0=\allowbreak 23/20,\allowbreak s_0=\allowbreak 0$, the linear coefficient increases to at most
$66/5$ on $[0,\allowbreak 2]$ and the merged quadratic is
$(13/5)r_0^2+(36/5)r_0<\allowbreak 12$.
Compare with $8e^{(79/16)t^2}\ge\allowbreak 8+(79/2)t^2$;
$8/5-(66/5)t+(55/2)t^2>\allowbreak 0$, with discriminant $-44/25$.
The signed covariance uses
$\operatorname{Re}(G_aQ)\ge\allowbreak 1-2x-(2-x)x^4-x\delta_3$.
After the true mean correction,
\fitdisplay{
\operatorname{Cov}(C_a,Y_1)
\le-(3/16-2\theta^4)d+[\theta^3-7\kappa/8]d^2
\le-13d/100-2d^2/5.
}{}
Here $\theta^4<\allowbreak 1/36$ gives $3/16-1/18=\allowbreak 19/144>\allowbreak 13/100$,
and $\theta^3<\allowbreak 7/100$.
The exact norm bound is
\fitdisplay{
|Q|^2\le1+x+4x^4+x^7+x^8
+2\delta_3x(1+x^3)+\delta_3^2x.
}{}
The linear coefficient is less than $61/40$, using
$\theta^4<\allowbreak 11/400$, $\theta^7<\allowbreak 1/500$, $\theta^8<\allowbreak 1/1000$.
The cross quadratic contribution is at most
$2\theta^2(\theta+\theta^4)\le\allowbreak 4/27$;
after subtracting $\zeta^2\ge\allowbreak 49d^2/64$, it is at most $-3d^2/5$.
The cubic contribution satisfies
$\delta_3^2\delta_2\le\allowbreak \theta^5d^3<\allowbreak d^3/80$.
Therefore $\operatorname{Var}Y_1\le\allowbreak  V_*:=\allowbreak 61d/40-3d^2/5+d^3/80$.
Use $\mu^3\ge\allowbreak 1-12d/5$ and
$\mu^4\le\allowbreak (1-d/2)^4\le\allowbreak 1-2d+3d^2/2$.
The mean-plus-covariance bracket is at most
$\overline A=\allowbreak 1-(1483/400)d+(3523/1000)d^2+(1647/800)d^3$.
Set $H=\allowbreak 1+5d+(25/2)d^2\le\allowbreak  e^{2\Gamma}$.
With coefficient eight,
$G=\allowbreak H^2-\overline A-8V_*H=\allowbreak d[g_1+g_2d+g_3d^2+g_4d^3+g_5d^4]$,
where
\fitdisplay{
g_1=603/400>3/2,\quad g_2=-9723/1000>-10,\quad
 g_3=-4527/800>-6,\quad g_4=863/4,\quad g_5=-5/4.
}{}
For $d\le\allowbreak 1/4$, $g_4+g_5d>\allowbreak 215$, so the bracket is at least
$P(d)=\allowbreak 3/2-10d-6d^2+215d^3$.
At $d\le\allowbreak 1/16$ it exceeds $1/2$.
At $d\ge\allowbreak 1/16$, $P''=\allowbreak -12+1290d>\allowbreak 64$.
For $d_0=\allowbreak 2/15$, $P(d_0)=\allowbreak 1/6-8/75+1720/3375$ and
$P'(d_0)=\allowbreak -2/15$; hence $P(d)\ge\allowbreak  P(d_0)-1/7200>\allowbreak 1/2$.
Thus $G\ge\allowbreak  d/2$. Since $H(1/4)=\allowbreak 97/32<\allowbreak 13/4$, $H^2<\allowbreak 11$,
the reference fourth is at most $1-d/22$, in particular at most $1-d/32$.
In all four unit cases, the mean-plus-covariance upper is nonnegative
before replacing $e^{2\Gamma}$ by $H$.
In part Q178.5.1, $A\ge\allowbreak 45/64$; in Q178.5.2, $A(d,\allowbreak v)>\allowbreak 0$ and
$\overline A\ge\allowbreak  A(d,\allowbreak \theta)$; in Q178.5.3,
$\overline A\ge\allowbreak 1-47d/16\ge\allowbreak 17/64$; in Q178.5.4,
$\overline A\ge\allowbreak 1-(1483/400)d\ge\allowbreak 117/1600$.
Each $V_*$ is nonnegative. Therefore the replacements of $\alpha^4$
by $H^{-2}$ and $\alpha^2$ by $H^{-1}$ preserve their directions.

\paragraph{Q178.6 The high domain and the complete local table}\label{part:Q178.6}
For $1/4\le\allowbreak  d\le\allowbreak 1/3$, $c\ge\allowbreak 5$ gives
$h_c=\allowbreak T f_{c-1}\le\allowbreak 2\E f_{c-1}\le\allowbreak 2(1-2d)\le\allowbreak 1$.
With the same canonical direction,
$T^L(qf_bh_c)\le\allowbreak  T^L(qf_b)$.
For the actual $a\ge\allowbreak3,\allowbreak b\ge\allowbreak3,\allowbreak L\ge\allowbreak1$,
\eqref{eq:q117-unified:shortone} gives
\fitdisplay{
\E[f_a\{T^L(qf_b)\}^4]\le1-d/3000.
}{}
Only that full $(1,\allowbreak b)$ interface is used; the short $(1,\allowbreak 2)$ or unit
$h_2f_b$ rows are not premises of this comparison.
The low-domain bounds are
\fitdisplay{
\begin{array}{c|c}
\text{original window}&0<d\le1/4\text{ reference fourth upper}\\\hline
b\ge4,c>b,L\ge2&1-d/16\\
L=1,b=4,c=5&1-d/16\\
L=1,b=4,c\ge6&1-d/160\\
L=1,b=5,c\ge6&1-d/2600\\
L=1,b\ge6,c>b&1-d/32.
\end{array}
}{}
The high domain gives $1-d/3000$ for every window.
At $d=\allowbreak 0$ only $K\one=\allowbreak 0$ is implied; original leaf integration makes
all actual messages and equal-pair products one.
Combining these domains proves \eqref{eq:Q178:T186.A} for all original
integers. The common positivity restoration of $E_t,\allowbreak E_o\in[0,\allowbreak 1]$
and the complete reversal of the original kernel, receiving arm,
terminal and connector give the decorated and reversed versions in the
same original finite measure.

\paragraph{The genuine receiving arm of length two}
\label{part:Q178:f2-extension}
For the same actual terminal \(F=qf_bh_c\), \(4\le b<c\),
and every \(L\ge1\),
\begin{equation}\label{eq:Q178:f2-extension}
 \E[f_2\{T^LF\}^4]\le1.
\end{equation}
On \(0<d\le1/4\), the proof gives the stronger common bound
\(1-d/3000\). On the high domain the assertion in
\eqref{eq:Q178:f2-extension} uses only the genuine length-two
reference bound by one. The strict all-domain \(d/3000\) loss
for receiving lengths \(a\ge3\) remains the assertion above.

Put \(C_2=f_2=1-u+Bu\), \(\E C_2=1-d\), and
\(C_2\le p+\sqrt d\). The coefficient templates in
paragraph~\ref{part:Q178.3.4} assume only
\(C_a\le p+(6/5)\sqrt d\). The coefficient eight in
paragraph~\ref{part:Q178.3.5} uses this same weak cap.
Thus all the actual scaled coefficients hold for \(C_2\),
without a positive lower bound on \(\beta_2\).
The packet \(r,R,D,\Gamma\), finite \(\gamma\)'s, positive
uppers, actual means, and variances remain those constructed
in paragraphs~\ref{part:Q178.2}, \ref{part:Q178.4}, and
\ref{part:Q178.5}. Only the receiving covariance and mass change.

For \(L\ge2\), the following unscaled covariance bounds follow
from \eqref{eq:reference-common:covariance}, with \(a=2,t=0\):
\begin{equation}\label{eq:Q178:f2-nonunit-budgets}
\begin{array}{c|c|c}
L&\text{new covariance upper}&\text{comparison budget }v_0\\\hline
2&-33d/800&29d/256\\
3&hd-d/4&hd+\rho\theta^{3/2}d\\
4&2hd-7d/25&d^2+d/20\\
\ge5&3hd-3d/10&3hd+d/30,
\end{array}
\qquad h=15d/32.
\end{equation}
The \(L=2\) comparison budget is the smallest of its three
actual \(b\) rows; the other two allow larger values.
To verify the new column, use the finite genuine products
\(\operatorname{Re}(G_2G_3)\ge3/16\) and
\(\operatorname{Re}(G_2G_n)\ge\kappa^2\) for \(n\ge4\).
The latter follows by the finite even/odd formulas as in
paragraph~\ref{part:Q175:f2-extension}.
At length two the full symbol \(G_3+D+z^2r\) has product at
least \(3/16-(6/5)(3/10)\theta=33/800\).
At length three, \(DG_2\) is favorable and
\(\kappa^2-(6/5)(3/10)(4/15)>1/4\).
At length four, \(DG_3\) is favorable and
\(\kappa^2-(6/5)(3/10)\theta^2>7/25\).
For all \(L\ge5\), \(DG_{L-1}\) is favorable and
\(\kappa^2-(6/5)(3/10)\theta^{5/2}>3/10\).
The retained shifts cost respectively \(h,2h,3h\) in the last
three rows. Each actual signed mean parameter is positive,
as proved for these same upper functions in paragraph~\ref{part:Q178.4}.

The comparison mass in every nonunit row is \(1-19d/16\),
and the actual means are at most one. Each improvement in
\eqref{eq:Q178:f2-nonunit-budgets} exceeds \(3d/64\).
Lemma~\ref{lem:reference-common:mass-transfer} therefore applies
before the actual mean substitutions. The same variances and
scaled coefficients give the identical scalar upper
\(e^{-25d/2}+(103/10)d e^{-9d/2}\le1-d/16\).

At \(L=1\), retain all four actual finite configurations:
\begin{equation}\label{eq:Q178:f2-unit-budgets}
\begin{array}{c|c|c}
(b,c)&\text{new covariance upper}&v_0\\\hline
(4,5)&-d/2&0\\
(4,c),\ c\ge6&-d/2-4d^2/5&-4d/25-2d^2/5\\
(5,c),\ c\ge6&-\kappa d-3d^2/4&-3d/16-3d^2/8\\
b\ge6&-53d/100-4d^2/5&-13d/100-2d^2/5.
\end{array}
\end{equation}
Write \(r=A_0+zC_0\). The full output symbol is
\(Q=(1-xC_0)+z(1+A_0)\), and its actual mean parameter is
\(\zeta=\int(1+A_0)\,d\nu>0\).
The first row has \(Q=1+z(1+x^2)\), hence product at least
\(1-x-x^3>1/2\).
For \(b=4,c\ge6\), the finite relation
\(U-(1-x)V\le x\) and \(|c_0|\le\delta_2\) give
\[
 \operatorname{Cov}(C_2,Y_1)
 \le-(\kappa-\theta^4)d-(1-\theta^2)d^2
 \le-d/2-4d^2/5.
\]
For \(b=5\), the negative finite coordinates give
\(A_0+C_0\le\delta_2\) and \(\zeta\ge11d/12\), so the upper
is \(-\kappa d-(11/12-\theta^2)d^2\).
For \(b\ge6\), use
\(|A_0|\le\delta_3+x^3\), \(|C_0|\le x^3\), and
\(\zeta\ge7d/8\). The upper is then
\[
 -(\kappa-2\theta^4)d-(7/8-\theta^3)d^2
 \le-53d/100-4d^2/5.
\]
These calculations retain the signed term \(-d\zeta\).

Every unit-row improvement exceeds \(d/4\). The comparison
mass is either the same-host \(1-2(1-v)d\), \(v=\delta_2/d\),
or its conservative endpoint \(1-19d/16\). Actual means are
at most one. The mass-transfer lemma therefore pays the full
change before any mean or \(v\) endpoint is taken.
The four scalar computations of paragraph~\ref{part:Q178.5},
with their actual negative variance squares and the stronger
\(\Gamma\ge(2601/1024)d>5d/2\) in the last row, now give
deficits \(d/16,d/160,d/2600,d/22\), respectively.
When clearing a scalar denominator, use the same nonnegative
comparison bracket established there; the scaled variance term
retains its second-power scale. This proves the stated common
low-domain deficit.

For \(d\ge1/4\), the same actual arm has \(h_c\le1\), so
\(T^L(qf_bh_c)\le T^L(qf_b)\). The genuine \(f_2\)-weighted
input \eqref{eq:L58:claim} gives a moment at most one for every
\(b\ge4,L\ge1\). At \(d=0\), use leaf integration.
Positivity restores the original equal-pair factors, and full
reversal reverses the receiving arm and connector together.
